\HMTNarrativeHeading{Una unidad, dos realizaciones}

La doble proyección organiza este artículo desde su primera construcción. Un
mismo estado conserva los datos que permiten precisar un valor y los datos
que permiten determinar una incidencia. La primera realización reúne
cilindros compatibles hasta individualizar su lectura numérica; la segunda
transporta las relaciones de frontera. La tesis estudia su procedencia
común y las operaciones que mantienen su correspondencia durante el
refinamiento. Número y configuración geométrica reciben así una genealogía
explícita.

La unidad se constituye mediante posiciones, operaciones y relaciones de
transporte. APP dispone las dos direcciones nonádicas sobre un mismo
soporte y evalúa sus hojas aditiva y multiplicativa. Cada evaluación
conserva el residuo y su cociente. TRIT distingue la orientación y el
régimen local; el TPK compone esos datos con el desplazamiento, el acarreo,
la frontera y la prolongación. La estructura discreta conjunta del
continuo reside en esa construcción enriquecida. Sus realizaciones
posteriores reciben operaciones ya determinadas.

La conexión nonádica tiene aquí una función constitutiva. Una vuelta
compone las transiciones del estado completo. La fase observable retorna
y el registro incorpora lo recorrido. Esta diferencia permite reconocer
una regularidad mientras la historia aumenta de profundidad. La
continuidad del refinamiento conserva los prefijos que hacen compatibles
las dos realizaciones. Cada precisión nueva pertenece a la misma historia
de la que proceden las anteriores.

El sentido de conservación evolutiva aparece en este punto. Una cantidad
normalizada puede permanecer igual mientras cambian la distribución, la
resolución o la representación. Recuperar el estado exige identificar qué
datos acompañan ese cambio. La aritmética conserva cocientes; el transporte
conserva orientación y vueltas; la incidencia conserva relaciones; el
archivo operatorio conserva los complementos necesarios para invertir la
lectura. El artículo desarrolla los mapas que articulan esas funciones.

La extrema y media razón holográfica expresa esta cuestión de unidad,
proporción y transformación. Su formulación incluye la completación
dimensional y el refinamiento reversible, y alcanza después los balances
de estados y observables. La constante de acoplamiento ocupa una posición
determinada dentro de ese recorrido: el registro generado vincula la
lectura aritmética, la incidencia y la recuperación operatoria. Las
demostraciones siguientes construyen los antecedentes de cada una de esas
funciones.

